\documentclass[twoside]{article}
\usepackage[preprint]{aistats2027}
\usepackage[round,authoryear]{natbib}
\usepackage{amsmath,amssymb,amsthm}
\usepackage{mathtools}
\usepackage{microtype}
\usepackage{booktabs}
\usepackage{url}
\renewcommand{\bibname}{References}
\renewcommand{\bibsection}{\subsubsection*{\bibname}}
\bibliographystyle{apalike}
\setlength{\bibsep}{0pt plus 0.3ex}
\newtheorem{theorem}{Theorem}
\newtheorem{proposition}[theorem]{Proposition}
\newtheorem{corollary}[theorem]{Corollary}
\newtheorem{lemma}[theorem]{Lemma}
\newtheorem{remark}[theorem]{Remark}
\newcommand{\Dens}{\mathcal D}
\newcommand{\M}{\mathcal M}
\newcommand{\C}{\mathcal C}
\newcommand{\R}{\mathbb R}
\newcommand{\tr}{\operatorname{tr}}
\newcommand{\supp}{\operatorname{supp}}
\newcommand{\KL}{\operatorname{KL}}
\newcommand{\Hsq}{\operatorname{H}^2}
\newcommand{\Ent}{\operatorname{S}}
\newcommand{\Ber}{\operatorname{Ber}}
\newcommand{\ip}[2]{\left\langle #1,#2\right\rangle}
\newcommand{\norm}[1]{\left\lVert #1\right\rVert}
\begin{document}
\runningauthor{}
\twocolumn[
\aistatstitle{Tight Logarithmic Bounds for Adversarial Online Phase Retrieval}
\aistatsauthor{Pahan Dewasurendra}
\aistatsaddress{Johns Hopkins University}
]
\begin{abstract}
We study realizable online phase retrieval: after an adversarial sensing vector $x_t$, a learner predicts the amplitude $|\ip{w_*}{x_t}|$ and then observes it exactly. Recent work proves an $O(\log^4 d)$ information-theoretic bound for bounded two-ReLU regression and an $\Omega(\log d)$ lower bound using the absolute-linear subclass. We close this gap for that subclass with an efficiently implementable horizon-free bound. Our algorithm lifts amplitudes to binary measurements of a density matrix and projects onto each revealed measurement constraint in quantum relative entropy. A quantum Pythagorean inequality and data processing make the measured Kullback--Leibler divergences telescope. Consequently, cumulative KL and squared Hellinger loss for arbitrary adaptive positive-operator-valued measurements are at most $D(\rho_*\|\rho_1)$. For a uniform prior this is $\log d-S(\rho_*)$. We prove that the KL value is exactly $\log(d/r)$ on maximally mixed rank-$r$ targets and give an effective-dimension Hellinger lower bound, making the entropy adaptation minimax sharp. The phase lift gives total squared amplitude loss at most $\log d$ for every horizon. A fixed target chosen before play and a nonanticipating dyadic sensing strategy force $\frac14\lfloor\log_2d\rfloor$ expected loss in $d-1$ rounds, with a matching lower-tail guarantee. Thus the worst-horizon minimax loss is $\Theta(\log d)$. We also prove a certified-interval upper bound of order $\log d+\sum_t\epsilon_t^2$ and a matching $\Omega(\log d+T\epsilon^2)$ lower bound for constant radius. Binary updates reduce to one-dimensional matrix-exponential root finding. The result requires exact probabilities or certified intervals, not sampled quantum outcomes.
\end{abstract}

\section{Introduction}

Phase retrieval asks for an unknown signal from magnitudes of linear measurements. Most theory assumes random sensing vectors and studies batch recovery or stochastic optimization \citep{candes2015phase,tan2019online,huang2023adversarial}. We ask a different question: if sensing vectors are chosen adaptively and amplitudes are revealed online, how much prediction loss is intrinsically necessary?

This question recently appeared inside realizable online regression. Every absolute-linear function has a two-ReLU representation
\begin{equation}
|\ip{w}{x}|=\operatorname{ReLU}(\ip{w}{x})+\operatorname{ReLU}(-\ip{w}{x}).
\label{eq:two-relu}
\end{equation}
For bounded $k$-ReLU networks, \citet{doronarad2026online} prove an $O(k^2\log^4 d)$ cumulative square-loss upper bound and an $\Omega(k^2\log(d/k^2))$ lower bound. Their hard fixed-width subclass is \eqref{eq:two-relu}. The upper bound is not computationally efficient and leaves three powers of $\log d$ unresolved even on this structured subclass.

The natural lift $|w^*x|^2=\tr(xx^*ww^*)$ turns phase retrieval into online prediction of a linear functional of a positive semidefinite matrix. This alone does not resolve the problem. Matrix exponentiated gradient already gives logarithmic realizable guarantees for \emph{squared intensity} error \citep{tsuda2005matrix}, while the desired amplitude loss applies a square root at zero and has no finite global Lipschitz constant. Generic contraction or squared-intensity arguments lose precisely where an amplitude can vanish.

Our key observation is that square-root geometry is native to probability distributions. Let a density matrix $\rho_*$ generate an exact outcome law $q_t$ under an adaptively chosen measurement. Starting from $\rho_t$, project in quantum relative entropy onto the affine set of states that generate $q_t$. The projection identity, measurement data processing, and the relation between Kullback--Leibler (KL) and Hellinger divergences give the one-step chain
\begin{equation}
\Hsq(p_t,q_t)\leq\KL(q_t\|p_t)
\leq D(\rho_*\|\rho_t)-D(\rho_*\|\rho_{t+1}).
\label{eq:intro-chain}
\end{equation}
It telescopes without a learning rate, horizon, spectral gap, or lower bound on the probabilities.

This yields four contributions. First, we prove the general KL and Hellinger telescope for arbitrary adaptive finite-outcome measurements, including singular iterates. The uniform-prior bound is $\log m-S(\rho_*)$, and we prove it exactly minimax for KL on every maximally mixed rank stratum under a fixed-target protocol. Second, a binary phase lift converts the first Hellinger coordinate exactly into amplitude error, giving an efficient $\log d$ bound for unit-norm phase retrieval and a $\log(d+1)$ bound when the signal norm is unknown. The argument extends to root-quadratic prediction. Third, we strengthen the dyadic construction of \citet{doronarad2026online}, specialized to absolute-linear functions, to a target chosen before play and a sensing strategy independent of learner predictions. It gives matching logarithmic expected and lower-tail bounds. This yields a tight worst-horizon characterization for the subclass that witnesses their general two-ReLU lower bound, not a new upper bound for all two-ReLU networks. Fourth, projection onto an amplitude interval gives a deterministic robust guarantee
\begin{equation}
\sum_t(\widehat y_t-y_t)^2
\leq\left(\sqrt{D(\rho_*\|\rho_1)}+2\sqrt{\sum_t\epsilon_t^2}\right)^2.
\label{eq:intro-robust}
\end{equation}
A combined construction proves that both dependencies are necessary up to constants.

The feedback model matters. Online quantum-state learning usually receives noisy estimates or sampled outcomes and controls mistakes or regret \citep{aaronson2018online,chen2020practical,yang2020revisiting,meyer2025pure}. Our finite telescope assumes that the target outcome distribution is revealed exactly, as an online regression label. Certified sets are covered by \eqref{eq:intro-robust}. A single sampled outcome does not define a consistency set containing the target and is outside our theorem.

\section{Entropic Outcome Prediction}
\label{sec:quantum}

Let $\Dens_m$ be the density matrices on $\mathbb C^m$, namely positive semidefinite Hermitian matrices of trace one. All results also hold over real symmetric matrices. A finite-outcome positive-operator-valued measure (POVM) is $\M=(E_1,\ldots,E_K)$ with $E_i\succeq0$ and $\sum_iE_i=I$. It maps a state to the distribution
\begin{equation}
\M(\rho)_i=\tr(E_i\rho).
\end{equation}
For distributions $p,q$, use the unnormalized squared Hellinger divergence
\begin{equation}
\Hsq(p,q)=\sum_i(\sqrt{p_i}-\sqrt{q_i})^2.
\label{eq:hellinger}
\end{equation}
For states $\rho,\sigma$, quantum relative entropy is
\begin{equation}
D(\rho\|\sigma)=\tr\{\rho(\log\rho-\log\sigma)\},
\label{eq:qre}
\end{equation}
with value $+\infty$ unless $\supp(\rho)\subseteq\supp(\sigma)$. The von Neumann entropy is $\Ent(\rho)=-\tr(\rho\log\rho)$.

\paragraph{Protocol.} Fix an unknown $\rho_*\in\Dens_m$. The learner initializes a full-rank $\rho_1$. On round $t$, an adaptive adversary reveals a POVM $\M_t$. The learner predicts $p_t=\M_t(\rho_t)$, observes the exact vector $q_t=\M_t(\rho_*)$, and incurs either $\KL(q_t\|p_t)$ or \eqref{eq:hellinger}. Define
\begin{equation}
\C_t=\{\rho\in\Dens_m:\M_t(\rho)=q_t\}.
\label{eq:consistency}
\end{equation}
The update is the information projection
\begin{equation}
\rho_{t+1}\in\arg\min_{\rho\in\C_t}D(\rho\|\rho_t).
\label{eq:iprojection}
\end{equation}
After a boundary outcome, minimization is on the support face of $\rho_t$. The target remains in that face, so the update is always well defined. Algorithmically, $\rho_t$ is chosen from past feedback before $\M_t$ is revealed. Its prediction after seeing $\M_t$ is proper and immediate.

\begin{theorem}[Entropic telescope]
\label{thm:telescope}
For every target, horizon, and adaptive sequence of POVMs, update \eqref{eq:iprojection} satisfies
\begin{align}
\sum_{t=1}^T\KL(q_t\|p_t)&\leq D(\rho_*\|\rho_1),
\label{eq:kl-telescope}\\
\sum_{t=1}^T\Hsq(p_t,q_t)&\leq D(\rho_*\|\rho_1).
\label{eq:h-telescope}
\end{align}
With $\rho_1=I/m$, both right sides equal $\log m-\Ent(\rho_*)\leq\log m$.
\end{theorem}

\begin{proof}
Quantum relative entropy is the Bregman divergence of negative von Neumann entropy. Projection onto \eqref{eq:consistency}, with the support-face convention, gives
\begin{equation}
D(\rho_*\|\rho_t)\geq D(\rho_*\|\rho_{t+1})+D(\rho_{t+1}\|\rho_t).
\label{eq:pythagoras}
\end{equation}
The POVM is a quantum-to-classical channel. Data processing and feasibility of $\rho_{t+1}$ yield
\begin{equation}
D(\rho_{t+1}\|\rho_t)\geq
\KL(\M_t(\rho_{t+1})\|\M_t(\rho_t))
=\KL(q_t\|p_t).
\label{eq:data-processing}
\end{equation}
Finally, if $B(p,q)=\sum_i\sqrt{p_iq_i}$, Jensen's inequality and $-\log z\geq1-z$ give
\begin{equation}
\KL(q\|p)\geq-2\log B(p,q)\geq2(1-B(p,q))=\Hsq(p,q).
\label{eq:kl-hellinger}
\end{equation}
Combining \eqref{eq:pythagoras}--\eqref{eq:kl-hellinger} proves \eqref{eq:intro-chain}. Summation telescopes. The supplement proves \eqref{eq:pythagoras} when the projection is singular and shows that zero predicted probabilities cause no undefined term.
\end{proof}

Theorem \ref{thm:telescope} is a realizable, comparator-dependent statement rather than regret to the best state after arbitrary labels. It is also stronger than a thresholded mistake bound: for every $\epsilon>0$, at most $D(\rho_*\|\rho_1)/\epsilon^2$ rounds can have Hellinger distance above $\epsilon$.

\section{Adversarial Phase Retrieval}
\label{sec:phase}

Fix $\mathbb F\in\{\R,\C\}$. On round $t$, an adversary reveals $x_t\in\mathbb F^d$ with $\norm{x_t}_2\leq1$. A fixed target $w_*\in\mathbb F^d$ generates
\begin{equation}
y_t=|\ip{w_*}{x_t}|.
\end{equation}
The learner predicts $\widehat y_t\in[0,1]$, observes $y_t$, and incurs $(\widehat y_t-y_t)^2$. Sensing vectors may depend on the entire previous transcript.

\begin{corollary}[Horizon-free phase retrieval]
\label{cor:phase-upper}
If $\norm{w_*}_2=1$, there is a lifted learner, efficiently implementable to certified numerical precision, such that, for every $T$,
\begin{equation}
\sum_{t=1}^T(\widehat y_t-y_t)^2\leq\log d.
\label{eq:phase-logd}
\end{equation}
If only $\norm{w_*}_2\leq1$ is known, the right side is at most $\log(d+1)$ and more precisely equals
\begin{equation}
\log(d+1)-h_2(\norm{w_*}_2^2),
\label{eq:norm-adaptive}
\end{equation}
where $h_2$ is binary entropy.
\end{corollary}

\begin{proof}
First suppose $\norm{w_*}_2=1$. Set $A_t=x_tx_t^*\preceq I$, use the binary POVM $(A_t,I-A_t)$, and take $\rho_*=w_*w_*^*\in\Dens_d$. Then its first target probability is
\begin{equation}
q_t=\tr(A_t\rho_*)=y_t^2.
\end{equation}
Run \eqref{eq:iprojection} from $I/d$ and predict
\begin{equation}
\widehat y_t=\sqrt{p_t},\qquad p_t=\tr(A_t\rho_t).
\end{equation}
The first coordinate of binary Hellinger divergence gives
\begin{equation}
(\widehat y_t-y_t)^2=(\sqrt{p_t}-\sqrt{q_t})^2
\leq\Hsq(\Ber(p_t),\Ber(q_t)).
\end{equation}
Theorem \ref{thm:telescope} and $\Ent(\rho_*)=0$ prove \eqref{eq:phase-logd}.

For $\norm{w_*}_2\leq1$, work in dimension $d+1$ with
\begin{equation}
\rho_*=\begin{pmatrix}w_*w_*^*&0\\0&1-\norm{w_*}_2^2\end{pmatrix},
\qquad
A_t=\begin{pmatrix}x_tx_t^*&0\\0&0\end{pmatrix}.
\label{eq:slack-lift}
\end{equation}
The two nonzero eigenvalues of $\rho_*$ are $\norm{w_*}_2^2$ and $1-\norm{w_*}_2^2$, which proves \eqref{eq:norm-adaptive}.
\end{proof}

The lift is not limited to rank one. If $Q_*\succeq0$ and $\tr(Q_*)=1$, the target $y_t=\sqrt{x_t^*Q_*x_t}$ has total squared loss at most $\log d-\Ent(Q_*)$. If $\tr(Q_*)\leq1$, the slack construction gives the analogous $d+1$ bound. Thus the theorem learns root-quadratic forms and adapts to their spectral entropy.

\section{A Matching Lower Bound}
\label{sec:lower}

Define the fixed-target worst-horizon minimax value of unit-norm phase retrieval by
\begin{align}
\mathfrak V_d(\mathbb F)
&=\sup_{T\geq1}\inf_{\mathcal A}\sup_{\substack{\norm{w_*}_2=1,\ w_*\text{ fixed}\\\text{adaptive sensing strategy}}}\nonumber\\
&\quad\mathbb E\!\left[\sum_{t=1}^T(\widehat y_t-y_t)^2\right].
\label{eq:minimax}
\end{align}
In the supremum, one target $w_*$ is fixed before play. The sensing vector $x_t$ may depend on past sensing vectors and revealed amplitudes, but not on learner predictions or private randomness. Its label is then $|\ip{w_*}{x_t}|$. The infimum includes randomized and improper learners.

\begin{theorem}[Fixed-target minimax and lower tail]
\label{thm:phase-minimax}
For every $d\geq2$ and $\mathbb F\in\{\R,\C\}$,
\begin{equation}
\frac14\lfloor\log_2d\rfloor\leq\mathfrak V_d(\mathbb F)\leq\log d.
\label{eq:minimax-sandwich}
\end{equation}
The lower bound uses at most $d-1$ rounds. More strongly, for every randomized learner and $\delta\in(0,1)$, there are a target fixed before play and a sensing strategy independent of its predictions for which, with probability at least $1-\delta$,
\begin{equation}
\sum_t(\widehat y_t-y_t)^2\geq\frac14\lfloor\log_2d\rfloor-\sqrt{\log\frac1\delta}.
\label{eq:phase-lower-tail}
\end{equation}
\end{theorem}

\begin{proof}
The upper bound is Corollary \ref{cor:phase-upper}. For the lower bounds, let $D=2^{\lfloor\log_2d\rfloor}$ and restrict to the first $D$ coordinates. Draw one target before play, uniformly from
\begin{equation}
w_s=\frac1{\sqrt D}(s_1,\ldots,s_D,0,\ldots,0),\qquad s\in\{-1,1\}^D.
\end{equation}
Use the dyadic merge invariant of \citet{doronarad2026online}. At any stage, partition $[D]$ into equal-size dyadic blocks. Each block $B$ has a public reference pattern $r_B\in\{-1,1\}^B$. Conditional on the labels so far, the target orientations on current blocks remain independent fair signs.

Merge two blocks $B_0,B_1$ of size $r/2$ by querying the unit vector
\begin{equation}
x_B=\frac1{\sqrt r}\left(\sum_{i\in B_0}r_{B_0}(i)e_i+
\sum_{i\in B_1}r_{B_1}(i)e_i\right).
\label{eq:merge-query}
\end{equation}
If the child orientations agree, the amplitude is $b=\sqrt{r/D}$. If they disagree, it is zero. Conditional on the past, these labels are fair and independent of the current prediction. For every prediction $a\in\R$,
\begin{equation}
\frac12\{a^2+(a-b)^2\}\geq\frac{b^2}{4}=\frac{r}{4D}.
\label{eq:merge-loss}
\end{equation}
The observed label records the corresponding relative orientation in the merged pattern, preserving conditional fairness across the new blocks.

There are $D/r$ merges at merged block size $r$, so each level contributes at least $1/4$ in conditional expectation. The $\log_2D$ levels use $D-1$ queries. Averaging over the fixed target gives one target attaining the expected lower bound against any randomized learner.

For \eqref{eq:phase-lower-tail}, clip each prediction to $[0,b]$, which only reduces its loss against $\{0,b\}$. If $L_t$ is its loss and $\mu_t$ its conditional mean, then $\mu_t\geq b^2/4$. The two possible values of $\mu_t-L_t$ lie in an interval of length at most $b^2$. Across all merges, $\sum_tb_t^4=2-2/D$. The conditional Hoeffding lemma therefore gives \eqref{eq:phase-lower-tail}. Averaging its success probability over the target prior again selects one fixed target. The supplement supplies the filtration and constants.
\end{proof}

Theorem \ref{thm:phase-minimax} closes the dimension dependence on the exact subclass used to prove the fixed-width lower bound of \citet{doronarad2026online}. Their construction chooses a maximizing label after each prediction. Conditional fairness strengthens it to a target fixed before play, a nonanticipating sensing strategy, and a lower tail. The result does not close the gap for arbitrary sums of two ReLUs, whose signs, offsets, and unequal directions do not reduce to one positive semidefinite rank-one target.

The phase construction also lower-bounds the general POVM game with pure real states and binary rank-one measurements. For general measurements, the entropy dependence admits an exact result. Let $\mathcal U_{m,r}$ be the states that are maximally mixed on an $r$-dimensional subspace, and define $\mathfrak W^{\KL}_{m,r}$ and $\mathfrak W^{\rm H}_{m,r}$ as in \eqref{eq:minimax}, with target in $\mathcal U_{m,r}$ and per-round loss respectively $\KL(q_t\|p_t)$ and $\Hsq(q_t,p_t)$.

\begin{theorem}[Exact KL value and entropy sharpness]
\label{thm:rank-minimax}
For $1\leq r\leq m$,
\begin{equation}
\begin{aligned}
\mathfrak W^{\KL}_{m,r}&=\log\frac mr,\\
(2-\sqrt2)\left\lfloor\log_2\left\lfloor\frac mr\right\rfloor\right\rfloor
&\leq\mathfrak W^{\rm H}_{m,r}\leq\log\frac mr.
\end{aligned}
\label{eq:rank-minimax}
\end{equation}
The KL lower bound uses one round.
\end{theorem}

\begin{proof}
Theorem \ref{thm:telescope} with $\rho_1=I/m$ gives both upper bounds because every target has entropy $\log r$. For the KL lower bound, draw $J$ uniformly from all rank-$r$ coordinate subsets before play and reveal the computational-basis POVM. For every prediction $p\in\Delta_m$,
\begin{equation}
\mathbb E_J\KL(\operatorname{Unif}(J)\|p)=-\log r-\frac1m\sum_{j=1}^m\log p_j\geq\log\frac mr
\end{equation}
by the geometric-mean inequality. Thus one fixed $J$ attains the lower bound. For Hellinger loss, put $k=\lfloor\log_2\lfloor m/r\rfloor\rfloor$, take $2^k$ disjoint coordinate blocks of size $r$, and draw the target block uniformly before play. Repeated balanced binary diagonal projectors make the deterministic target outcome conditionally fair. Every Bernoulli prediction has conditional expected loss at least $\min_p\{2-\sqrt p-\sqrt{1-p}\}=2-\sqrt2$. Summing and averaging over the fixed block proves the claim.
\end{proof}

\section{Intervals and Computation}
\label{sec:robust}

Exact amplitudes are idealized. Suppose instead that after predicting, the learner observes $z_t$ and a certified radius $\epsilon_t$ satisfying $|z_t-y_t|\leq\epsilon_t$. Let
\begin{equation}
I_t=[(z_t-\epsilon_t)_+,\min\{1,z_t+\epsilon_t\}]
\end{equation}
and project $\rho_t$ onto the linear slab
\begin{equation}
\C_t^{I}=\{\rho\in\Dens_m:\sqrt{\tr(A_t\rho)}\in I_t\}.
\label{eq:slab}
\end{equation}
This is convex because it equivalently constrains $\tr(A_t\rho)$ to the squared interval.

\begin{theorem}[Certified interval feedback]
\label{thm:robust}
For phase retrieval or root-quadratic prediction, the slab-projection learner satisfies
\begin{align}
\sum_{t=1}^T(\widehat y_t-y_t)^2
&\leq\left(\sqrt{D(\rho_*\|\rho_1)}+2\sqrt{\sum_{t=1}^T\epsilon_t^2}\right)^2\nonumber\\
&\leq2D(\rho_*\|\rho_1)+8\sum_{t=1}^T\epsilon_t^2.
\label{eq:robust-bound}
\end{align}
In particular, the unit-norm uniform-prior bound is $(\sqrt{\log d}+2\sqrt{\sum_t\epsilon_t^2})^2$, which reduces to $\log d$ under exact feedback.
\end{theorem}

\begin{proof}
Let $r_t=\tr(A_t\rho_{t+1})$. The target is in \eqref{eq:slab}, so projection and data processing give
\begin{equation}
\sum_t\Hsq(\Ber(p_t),\Ber(r_t))\leq D(\rho_*\|\rho_1).
\end{equation}
Both $\sqrt{r_t}$ and $y_t$ belong to $I_t$, hence $|\sqrt{r_t}-y_t|\leq2\epsilon_t$. By the triangle inequality in $\mathbb R^T$ and the first Hellinger coordinate,
\begin{align}
\left\|(\sqrt{p_t}-y_t)_{t=1}^T\right\|_2
&\leq\left\|(\sqrt{p_t}-\sqrt{r_t})_{t=1}^T\right\|_2\nonumber\\
&\quad+\left\|(\sqrt{r_t}-y_t)_{t=1}^T\right\|_2\\
&\leq\sqrt{D(\rho_*\|\rho_1)}+2\sqrt{\sum_t\epsilon_t^2}.
\end{align}
Squaring proves the first bound. The second follows from $(a+b)^2\leq2a^2+2b^2$.
\end{proof}

The dependence on uncertainty is minimax sharp. The following lower bound uses one target fixed before play and queries independent of learner predictions.

\begin{theorem}[Sharpness of certified intervals]
\label{thm:robust-lower}
Fix $d\geq4$, let $D=2^{\lfloor\log_2(d-2)\rfloor}$, and fix $N\geq1$ and $0<\epsilon\leq1/8$. Every randomized learner has a unit-norm target and a nonanticipating sequence of $D-1$ exact rounds followed by $N$ radius-$\epsilon$ rounds such that
\begin{equation}
\mathbb E\sum_t(\widehat y_t-y_t)^2\geq\frac18\log_2D+N\epsilon^2.
\label{eq:robust-lower}
\end{equation}
Hence the robust minimax dependence is $\Theta(\log d+N\epsilon^2)$ up to universal constants.
\end{theorem}

\begin{proof}
Draw signs $s\in\{-1,1\}^D$ and a branch $\nu\in\{-1,1\}$ independently and uniformly. Put squared norm $1/2$ on the sign block, amplitude $a+\nu\epsilon$ on coordinate $D+1$ for $a=1/4$, and the remaining norm on an unqueried coordinate. The target is unit norm because $1/2+(a+\epsilon)^2<1$. Scaled dyadic queries force expected loss $(1/8)\log_2D$. Then query $e_{D+1}$ for $N$ rounds and always return $z_t=a$ with certified radius $\epsilon$. The feedback is identical under both branches, and
\begin{equation}
\frac12\sum_{\nu\in\{-1,1\}}(\widehat y-a-\nu\epsilon)^2=(\widehat y-a)^2+\epsilon^2.
\end{equation}
Thus these rounds force $N\epsilon^2$ in expectation. Averaging over $(s,\nu)$ selects one fixed target. The supplement gives the combined construction in full.
\end{proof}

For a binary equality $\tr(A\rho)=q$ and an iterate positive definite on its support, the projection has the form
\begin{equation}
\rho(\lambda)=\frac{\exp(\log\rho_t-\lambda A)}{\tr\exp(\log\rho_t-\lambda A)}.
\label{eq:exponential-form}
\end{equation}
The function $\lambda\mapsto\tr(A\rho(\lambda))$ is nonincreasing because it is the negative derivative of a convex log-partition function. Bisection therefore reduces each phase update to a one-dimensional root problem, with an eigendecomposition or matrix exponential at each step. A slab update is a no-op if $p_t$ already lies in the interval, and otherwise projects to the nearer active boundary in the relative-entropy geometry. Boundary labels are limits on a reduced support face, as in minimum-relative-entropy quantum estimation \citep{zorzi2014minimum}. The supplement also gives an additive-defect certificate for approximate solves.

\section{Relation to Prior Work}
\label{sec:related}

\paragraph{Matrix entropic updates.} \citet{tsuda2005matrix} introduced matrix exponentiated-gradient updates, von Neumann Bregman projections, and a realizable logarithmic relative-loss analysis for squared errors in linear intensities. \citet{zorzi2014minimum} establish existence, uniqueness, support reduction, and exponential-family solutions for offline quantum relative-entropy estimation under affine constraints. \citet{sreekumar2026quantum} prove the finite-dimensional quantum Pythagorean theorem for compact convex and mixture families, including support conditions, in a batch maximum-likelihood setting. Our projection step uses this geometry. The new consequence is the online data-processed KL telescope, its multioutcome Hellinger form, and the square-root geometry that yields a tight amplitude result. Squared intensity control alone does not imply our horizon-free amplitude bound near zero.

\paragraph{Online quantum-state learning.} \citet{aaronson2018online} obtain $O((\log m)/\epsilon^2)$ large-error mistakes and $O(\sqrt{T\log m})$ generic regret for binary expectation prediction. Subsequent work gives practical, low-rank, pure-state, adaptive, or structured generic-regret refinements \citep{chen2020practical,yang2020revisiting,meyer2025pure,bansal2026structured}. \citet{bansal2024panoply} extend $O(\sqrt T)$ Lipschitz-loss regret to compact convex classes of positive semidefinite objects, including pure-state Gram matrices. \citet{tseng2025logarithmic} study the different loss $-\log\tr(A_t\rho_t)$ against the best hindsight state and obtain dimension-polynomial logarithmic regret. These protocols reveal noisy scalar outcomes or loss functions. None of their bounds states the finite cumulative KL or Hellinger guarantee under exact realizable outcome-law feedback. Conversely, our theorem does not improve their sampled-outcome guarantees.

\paragraph{Realizable regression and phase retrieval.} The scaled online dimension characterizes horizon-free realizable regression abstractly, and entropy integrals can bound it \citep{raman2025unified,doronarad2026online}. For bounded two-ReLU networks, the resulting dimension dependence is polylogarithmic. Our lifted information projection exploits the stronger algebraic structure of absolute-linear prediction to attain the exact logarithmic order efficiently. Stochastic phase-retrieval algorithms use random designs and convergence of a nonconvex iterate \citep{tan2019online}. \citet{bastianello2022opreg} use online phase retrieval as a time-varying optimization benchmark, tracking changing signals from noisy batch objectives rather than predicting adaptively chosen measurements. Robust batch work permits adversarial corruptions among random measurements \citep{huang2023adversarial,buna2025robust}. None of these settings permits adaptive worst-case sensing vectors with a horizon-free cumulative prediction guarantee.

\section{Discussion}

Relative entropy supplies exactly one logarithmic unit of initial uncertainty per effective dimension, while a dyadic adversary can extract a constant amount at each binary scale. This explains both sides of Theorem \ref{thm:phase-minimax}. The proof is insensitive to vanishing amplitudes because Hellinger distance measures square-root probabilities directly.

There are three important limits. First, exact probability feedback is stronger than observing one quantum measurement outcome. Second, the state lift is generally improper for phase retrieval because $\rho_t$ need not be rank one, though its amplitude prediction is efficiently computable. Third, the result resolves the absolute-linear hard subclass rather than the whole two-ReLU class. Extending the entropic telescope to signed combinations of several independently oriented rank-one components is a concrete obstacle left open.

\bibliography{references}

\clearpage
\appendix
\end{document}
